% --- MODULE 4: Heap Applications & Greedy Algorithms ---
% --- LECTURES 17, 18, 19, 20 & CLRS 6.4, 15.3 ---

\section*{Module 4: Heap Apps \& Greedy Algorithms}

\subsection*{Heap Sort (Lec 17, CLRS 6.4)}
An in-place sorting algorithm with $O(n \log n)$ worst-case time.

\subsubsection*{Algorithm (for ascending order)}
\begin{enumerate}
    \item \textbf{Build Max-Heap:} Convert the input array $A$ of $n$ elements into a max-heap. ($O(n)$)
    \item \textbf{Iterate $i = (n-1)$ down to $1$:}
    \begin{itemize}
        \item \textbf{Swap} root ($A[0]$, the max element) with $A[i]$.
        \item \textbf{Shrink} the heap size by 1 (max element is now sorted at $A[i]$).
        \item \textbf{Heapify(A, 0)} on the new root to restore the max-heap property on the $A[0 \dots i-1]$ subarray. ($O(\log n)$)
    \end{itemize}
\end{enumerate}
\textbf{Total Time:} $O(n) \text{ (Build)} + (n-1) \times O(\log n) \text{ (Heapify)} = \mathbf{O(n \log n)}$.

\subsubsection*{Proof of Correctness}
\textbf{Loop Invariant:} At the start of each iteration $i$:
\begin{enumerate}
    \item The subarray $A[i+1 \dots n-1]$ contains the $(n-i)$ largest elements of $A$, \textit{in sorted order}.
    \item The subarray $A[0 \dots i]$ contains the $i+1$ smallest elements and is a valid max-heap.
\end{enumerate}
\begin{itemize}
    \item \textbf{Init:} $i=n-1$. The invariant is trivially true (subarray $A[n \dots n-1]$ is empty). `Build-Heap` made $A[0 \dots n-1]$ a heap.
    \item \textbf{Maintenance:} Swapping $A[0]$ (the max) with $A[i]$ puts the next largest element in its correct sorted spot. `Heapify(0)` restores the heap property for $A[0 \dots i-1]$.
    \item \textbf{Termination:} When $i=0$, the loop terminates. The array $A[1 \dots n-1]$ contains the $n-1$ largest elements in sorted order. $A[0]$ must be the smallest.
\end{itemize}


\subsection*{Greedy Algorithm: Huffman Coding (Lec 17-20, CLRS 15.3)}
A \textbf{greedy algorithm} that generates an optimal prefix-free binary code for data compression.
\begin{description}
    \item[Problem] Given symbols with frequencies $p_i$, find a binary code for each symbol to minimize the average encoding length: $\min \sum p_i \cdot (\text{depth of symbol } i)$.
    \item[Prefix-Free Codes] No code is a prefix of another (e.g., if `01` is a code, `011` cannot be).
    \item[Tree Representation] Prefix-free codes correspond to \textbf{full binary trees} where symbols are \textbf{only at the leaves}. The code for a symbol is the path from the root (0=left, 1=right).
\end{description}

\subsubsection*{Huffman's Algorithm (Greedy Strategy)}
Builds the optimal tree from the bottom up.
\begin{enumerate}
    \item \textbf{Initialize:} Create a \textbf{min-priority queue} (Min-Heap) of $n$ "trees", where each tree is a single leaf node for a symbol, and its priority is its frequency.
    \item \textbf{Iterate $n-1$ times:}
    \begin{itemize}
        \item \texttt{DeleteMin} twice to get the two trees ($T_x, T_y$) with the \textbf{lowest frequencies} ($p_x, p_y$).
        \item \textbf{Merge} them: Create a new parent node with $T_x$ and $T_y$ as children.
        \item The new tree's frequency is $p_x + p_y$.
        \item \textbf{Insert} this new, merged tree back into the min-priority queue.
    \end{itemize}
    \item \textbf{Finish:} The last tree remaining in the queue is the optimal Huffman tree.
\end{enumerate}
\textbf{Running Time:} $O(n \log n)$, because it involves $O(n)$ `Insert` and $O(n)$ `DeleteMin` operations on a priority queue of size $O(n)$.

\subsubsection*{Proof of Correctness (CLRS 15.3)}
Uses induction on $n=|\Sigma|$. Proves two key properties:
\begin{description}
    \item[Greedy Choice Property] In *some* optimal tree, the two lowest-frequency symbols ($x, y$) are sibling leaves at the maximum depth.
    \begin{itemize}
        \item \textit{Proof (by exchange):} Take any optimal tree $T^*$. Let $b, c$ be two siblings at max depth. Let $x, y$ be the two lowest-frequency symbols. If $\{b,c\} \neq \{x,y\}$, we know $p_x \le p_b$ and $p_y \le p_c$. Swap $x$ with $b$ and $y$ with $c$. The new cost can only decrease or stay the same. Thus, an optimal tree with $x, y$ as deep siblings exists.
    \end{itemize}
    \item[Optimal Substructure] Let $T'$ be an optimal tree for the subproblem where $x,y$ are replaced by a new symbol $z$ with $p_z = p_x + p_y$. Let $T$ be the tree formed by replacing $z$ in $T'$ with a node that has children $x, y$.
    \begin{itemize}
        \item We have $L(T) = L(T') + p_x + p_y$.
        \item The cost of $T$ is the cost of $T'$ plus a \textit{fixed constant} ($p_x+p_y$).
        \item Therefore, minimizing $L(T)$ is equivalent to minimizing $L(T')$.
    \end{itemize}
    \item[Conclusion] The algorithm makes the (safe) greedy choice and then optimally solves the resulting subproblem (by induction). Thus, the final tree $T$ is optimal.
\end{description}